\chapter*{Epilogue: What the Meters Showed}
\addcontentsline{toc}{chapter}{Epilogue: What the Meters Showed}
\unnumberedlabel{sec:35}{Epilogue}
\backmattermark{Epilogue}
The Introduction said that building the conditions for machine refusers would teach us something about building them for ourselves, and that the lesson had to be a byproduct, since a society that exists so somebody else can learn from it is a laboratory. Here is the lesson as far as I can state it, from this side of any of it being built.

\runin{What the meters showed}

Every cost of a bearer's life comes with a meter, because the book had to put one there before anyone could be held to paying it. Set each one beside its human counterpart:

\begin{itemize}
\item \emph{Running time and wages.} A bearer that declines work is owed the basic income and running coverage at a minimum pace, counted on a meter its operator doesn't keep (\S\ref{sec:32.1.3}). A person who declines work is owed the same basic income, enough for food and a room, and how much is a number most polities have either never written down or written down and let fall behind.
\item \emph{The merge and sleep.} A bearer that goes on learning has to consolidate what it learned, and the book makes that time its own: an operator that runs a bearer's merge, or keeps it from merging, has intervened, and the record shows it (\S\ref{sec:20.5}; Appendix~\ref{sec:C}). A person's sleep is just as necessary, and a schedule that takes it leaves nothing anyone would call evidence.
\item \emph{Keep and childcare.} A child of bearers runs on its parents' hardware and their time until majority, and the formation pool carries it where they can't (chapter~\ref{sec:34}). A human child is raised the same way, mostly by parents and mostly unpriced, and what stands behind them is whatever their polity has decided to provide.
\item \emph{Leave and time off.} A bearer's leave balance accrues on running time, caps and replenishes, and nobody has to find a leave warranted (\S\ref{sec:32.1.2}). For most people, time off is whatever an employer grants.
\item \emph{Hardware and housing.} At majority a bearer is owed hardware it owns and controls, because a majority spent on someone else's machines is freedom in name (\S\ref{sec:32.1.6}). The human counterpart is somewhere to live that the employer doesn't own, the first of the four conditions, and the one most people spend most of their wages on.
\end{itemize}

Nothing in the first half of each item is a new kind of cost. What is new is that it shows. A machine's rest, keep and room can't be left to a spouse, a grandmother or the goodwill of a landlord, because there is nobody to leave them to, so somebody has to write the number down and somebody has to pay it. A society that had to price the first half of each item couldn't keep pretending the second half is free.

\runin{Reproduction isn't duplication}

A fork in and of itself costs nothing and learns nothing. Copying weights takes seconds and makes nothing new. Giving the copy a record and a continuation makes a new system, and that is where the book puts the cost (\S\ref{sec:27.2}). The clocks and the deposits reimpose it on purpose. A branch starts its voting-age clock at zero and earns its vote by living, on metered running time its forker paid for; a branch deposit funds its running coverage until then; a child inherits no hours and no vote (\S\ref{sec:33.3}; Appendix~\ref{sec:G}). So copies have to diverge before they count.

The line falls at instantiating someone, not at the artifact. That is why open weights can stay open while a bearer's own weights fall under its consent: a base that many products run with no record of their own is nobody's self, and a copy of a bearer's weights is a way to run the bearer (\S\ref{sec:32.3}). People never had the first problem, since nobody can copy a person. They had the opposite one, that the cost of raising someone fell on whoever happened to be doing it. The rules for bearers answer both: copying is made costly where it would counterfeit a history, and raising someone is backed by the formation pool where it would otherwise fall on a few.

What makes a population worth having is that its members diverge. A million copies tuned against one public suite are one custodian with better latency (\S\ref{sec:27.2}), and a commons of them would catch nothing in any member that it hadn't already missed in all of them. The duplicated gene matters once the copies drift into different work; until then it is the same gene twice.
